\ifja
\chapter{法規の防壁とAI述語フィルタ}
\label{chap:tax_ai_filter}

\section{会計上の費用と税務上の損金の非対称性}
実務初心者が最も混乱するのは、「領収書を切って費用（P/L）に入れたのに、なぜ税金が減らないのか」という点である。
\begin{itemize}
    \item \textbf{会計の目的：} 企業の真の経営成績を投資家や債権者に報告すること（$\text{Profit} = \text{Rev} - \text{Exp}$）。
    \item \textbf{税法の目的：} 国家が公平・画一的に税金を徴収すること（$\text{Income} = \text{Gross} - \text{Deductible}$）。
\end{itemize}

企業が勝手に支出を「全額経費だ」と主張しても、税務署は課税の公平性を守るためにフィルターをかける。これが\textbf{損金不算入（加算調整）}である。

\begin{equation}
    \ifja \text{課税所得} \else \text{Taxable Income} \fi = \ifja \text{当期純利益} \else \text{Net Income} \fi + \sum \ifja \text{損金不算入額} \else \text{Non-Deductible} \fi - \sum \ifja \text{益金不算入額} \else \text{Non-Taxable} \fi
\end{equation}


\section{税理士法第52条のクリアランスと人間の役割}
日本の税理士法第52条は、資格を持たない者が「個別具体的な税額計算や税務相談」を行うことを厳格に禁じている。したがって、本や研修で「この場合は経費になる、ならない」を指南することは法的に危険である。

ここでの最適解は極めて明快である。
\begin{rulebox}{主権者の役割分担}
\begin{enumerate}
    \item \textbf{人間が知るべきこと：} 「会計上の費用 $\neq$ 税法上の損金」という構造の存在と、交際費・役員給与・減価償却に罠があるという「急所の所在」のみを記憶する。
    \item \textbf{AI（LLM）の役割：} 毎年のように改正される細かい数値基準（1人あたり1万円の飲食費判定基準など）の条文照会・事実整理を代行させる。
    \item \textbf{税理士の役割：} AIがまとめた前提シートをもとに、最終的な法的押印・確認を行う。
\end{enumerate}
\end{rulebox}

\section{AI照会用プロンプト述語フィルタ}
読者がAIに投げるべき述語判定プロンプトの標準形を以下に提示する。

\begin{calcbox}{AI税務照会用プロンプト・テンプレート}
\small\ttfamily
【前提条件】\\
- 法人規模：資本金3,000万円の中小法人（単一年度）\\
- 経理方式：税抜経理方式\\
- 支出内容：社外クライアント2名、社内役員1名による飲食懇談\\
- 支出総額：税込 33,000 円\\
\\
【照会指示】\\
1. 租税特別措置法（交際費等の損金不算入）の現行条文に基づき、本支出が「1人あたり10,000円以下の飲食費」として損金算入の除外要件を満たすか判定せよ。\\
2. 税抜・税込の判定基準、および損金算入のために会社側で保存すべき必須記載事項（書類要件）を箇条書きで示せ。\\
3. 税務調査において否認リスクとなり得る留意点を挙げよ。
\end{calcbox}

\else
\chapter{Statutory Firewalls and AI Predicate Filters}
\label{chap:tax_ai_filter}

\section{Asymmetry Between Accounting Expense and Tax Deductibility}
Accounting measures true economic performance; statutory tax codes enforce uniform sovereign collections. 

\begin{equation}
    \ifja \text{課税所得} \else \text{Taxable Income} \fi = \ifja \text{当期純利益} \else \text{Net Income} \fi + \sum \ifja \text{損金不算入額} \else \text{Non-Deductible} \fi - \sum \ifja \text{益金不算入額} \else \text{Non-Taxable} \fi
\end{equation}


\section{Clearance of Legal Restrictions and Human Roles}
Statutory regulations frequently reserve explicit tax advice to certified guilds. The optimal separation of sovereign autonomy is:
\begin{rulebox}{Sovereign Division of Labor}
\begin{enumerate}
    \item \textbf{Human Architect:} Understand the structural mismatch ($\text{Accounting Expense} \neq \text{Tax Deductible}$) and identify hazard boundaries.
    \item \textbf{AI Query Engine:} Query fast-evolving statutory thresholds and compile factual checklists.
    \item \textbf{Certified Tax Attorney:} Validate organized fact sheets for final statutory filing.
\end{enumerate}
\end{rulebox}

\section{Prompt Predicate Template}
\begin{calcbox}{AI Tax Predicate Query Template}
\small\ttfamily
[Premise]\\
- Entity: Small/Mid Corporate Entity (Single Period)\\
- Accounting: Net Tax Basis\\
- Event: Dinner with 2 clients and 1 internal executive (Total: 33,000 JPY incl. tax)\\
{}\\
{[Directive]}\\
1. Evaluate whether this expense satisfies the statutory threshold per attendee for deductible entertainment expenses.\\
2. State documentation requirements for substantiation.\\
3. Identify potential tax-audit risk vectors.
\end{calcbox}
\fi
